\documentclass[sigconf, 10pt, anonymous, nonacm]{acmart}
\settopmatter{printfolios=true, printccs=false, printacmref=false}

\title{When Measurements Become Targets: The Strategic Manipulation Challenge for Internet Observatories}

\author{First Author}
\affiliation{%
  \institution{First Institution}
  \city{City}
  \country{Country}
}

\author{Second Author}
\affiliation{%
  \institution{Second Institution}
  \city{City}
  \country{Country}
}

\begin{abstract}
Internet observatories are vulnerable to strategic behavior because they rely on a publicly known and relatively static measurement infrastructure. As observatory outputs increasingly influence reputation, compliance, and business relationships, operators gain incentives to selectively optimize their behavior toward measurement platforms. This raises a fundamental question: are Internet observatories measuring network properties, or merely measuring how networks behave when they know they are being observed?
\end{abstract}

\begin{document}

\maketitle

\section{Introduction}

Internet measurement plays a fundamental role in understanding the operational state of the Internet. Researchers, network operators, regulators, and industry stakeholders rely on measurement platforms to assess a wide range of properties, including network performance (e.g., latency, packet loss, throughput), routing security (e.g., RPKI deployment, Route Origin Validation (ROV), Source Address Validation (SAV)), censorship and Internet freedom, resilience, and reliability. Over the past decade, numerous large-scale measurement efforts have been developed to provide visibility into these aspects of Internet behavior, including performance measurement platforms, routing security observatories, and censorship monitoring frameworks.
A common requirement across these efforts is the availability of geographically and topologically diverse measurement vantage points. Ideally, such probes should satisfy three properties. First, they should be freely or economically accessible to researchers. Second, they should provide broad geographic and network diversity to maximize Internet coverage. Third, they should allow controlled and reproducible experiments. These requirements significantly limit the available measurement infrastructures, leading researchers to rely heavily on a relatively small set of platforms such as RIPE Atlas, CAIDA Ark, M-Lab, and similar systems.
Alternative measurement techniques based on side channels, such as IP-ID counters or IPv6 rate-limiting behaviors, have also been proposed to expand visibility beyond dedicated measurement infrastructures. However, such techniques typically suffer from limited controllability, reduced accuracy, and declining availability as network implementations evolve. Furthermore, many Internet studies require longitudinal measurements spanning months or years, making stability of measurement vantage points a critical requirement. For example, long-term censorship studies often maintain a stable set of probes and replace only inactive nodes to preserve measurement consistency and reproducibility across time.
The outputs of these measurement platforms increasingly influence operational and business decisions. For instance, the MANRS Observatory derives several routing security metrics from publicly observable measurements and uses them to estimate the security posture of Autonomous Systems (ASes). These scores are intended to provide objective indicators of operational best practices and are frequently used by network operators, customers, regulators, and researchers to evaluate the trustworthiness of networks.
However, this growing reliance on external measurements introduces a fundamental assumption: networks behave identically toward measurement probes and ordinary users. Existing measurement methodologies implicitly assume that observed behavior reflects the network's actual operational policies. In this paper, we argue that this assumption is increasingly questionable.
Unlike ordinary users, many Internet measurement infrastructures are publicly documented. The IP addresses of probes, their geographic locations, network affiliations, and measurement activities are often known or can be inferred with high confidence. Consequently, network operators may be able to identify when they are being measured and selectively alter their behavior toward measurement traffic. This creates an opportunity for strategic manipulation of Internet observatories.
Consider Source Address Validation (SAV) as an example. The Spoofer Project measures the Internet's susceptibility to IP spoofing using a set of known measurement destinations and vantage points. These measurements are subsequently used by initiatives such as the MANRS Observatory to assess SAV deployment. A network operator seeking a favorable security score could simply identify the destinations associated with the measurement infrastructure and block spoofed traffic only toward those destinations, while continuing to permit spoofed traffic toward the rest of the Internet. Such a network would appear compliant to the measurement framework despite failing to deploy SAV comprehensively.
A similar situation can arise in censorship measurements. Suppose a censoring network is aware of the IP addresses used by a measurement platform such as RIPE Atlas. The network could selectively whitelist traffic originating from these probes while continuing to block access for ordinary users. Measurement platforms would conclude that the content remains accessible, whereas users within the network would experience censorship.
Performance measurements are equally vulnerable. A network operator could identify traffic originating from known measurement probes and prioritize it using dedicated Quality-of-Service (QoS) policies. Consequently, latency, packet loss, and throughput measurements would reflect artificially improved performance that does not represent the experience of ordinary users.
These examples share a common characteristic: the measurement infrastructure itself becomes part of the attack surface. Because measurement probes are relatively stable, publicly visible, and repeatedly reused across experiments, operators can potentially identify them and provide special treatment. Moreover, measurement campaigns often utilize only a subset of available probes at a given time, making detection even easier. If incentives exist—such as improving public reputation, increasing compliance scores, attracting customers, or satisfying regulatory requirements—operators may have strong motivations to behave strategically.
The challenge extends beyond any single measurement domain. Whether the objective is measuring routing security, censorship, performance, resilience, or operational best practices, most Internet observatories infer internal network behavior from externally observable outcomes. As a result, they are inherently vulnerable to strategic behavior by the entities they attempt to measure.
In this paper, we argue that Internet measurement frameworks should not only account for measurement errors and technical limitations, but also consider the possibility of adversarial or strategic behavior by measured networks. We introduce the notion of measurement-aware networks, in which operators can recognize measurement infrastructures and selectively adapt their behavior to influence observed outcomes. We show that such behavior can fundamentally distort conclusions drawn from Internet observatories and compliance frameworks, leading to inaccurate assessments of security, performance, censorship, and operational trustworthiness.
Ultimately, our work raises a broader question: How trustworthy are Internet measurements when the measured networks know they are being measured? Addressing this question is essential for ensuring that future Internet observatories remain robust against strategic manipulation and continue to provide accurate representations of Internet behavior.


This sentence shows an example citation~\cite{cerf-kahn-1974}. Each reference
must list all authors of the cited work.

\section{Second Section}

Replace this text with the body of your paper.

\section{Third Section}

Replace this text with the body of your paper.

\bibliographystyle{ACM-Reference-Format}
\bibliography{hotnets26-template}

\end{document}
